% Lösungen Flächeninhalt und Volumen im Raum · Zusammengesetzt und Verhältnisse · Prüfungs-Fokus Abitur GK (zusammenbau v0.6)
\einheitenkopf{Einheit 4 · Einheit 4}

\erg{1}{a) als Summe: Quader plus Dachprisma \quad b) Pyramide $PQUT$: $V = \tfrac{1}{3} \cdot 18 \cdot 8 = 48$; über der Deckfläche ragt die ähnliche Pyramide mit den Kanten 3, 3 und 4 heraus: $V = \tfrac{1}{3} \cdot 4{,}5 \cdot 4 = 6$; Teilkörper $= 48 - 6 = 42$ \quad c) In der Ebene gilt $z = \tfrac{y}{2}$; der kleinere Teilkörper ist ein Prisma der Länge 8 mit dreieckigem Querschnitt (Katheten 8 und 4): $V = \tfrac{1}{2} \cdot 8 \cdot 4 \cdot 8 = 128$}
\erg{2}{a) Pyramide $PQUT$: $V = \tfrac{1}{3} \cdot 18 \cdot 10 = 60$; über der Deckfläche ragt die ähnliche Pyramide mit den Kanten 3, 3 und 5 heraus: $V = \tfrac{1}{3} \cdot 4{,}5 \cdot 5 = 7{,}5$; Teilkörper $= 60 - 7{,}5 = 52{,}5$ \quad b) Die verlängerten Seitenkanten treffen sich in der Spitze der ganzen Pyramide; die Deckfläche halbiert die Seite, also liegt die Spitze 6\,m über der Deckfläche und 12\,m über der Grundfläche. Erster Summand: ganze Pyramide (Grundfläche $36\,\mathrm{m}^2$, Höhe 12\,m), zweiter: abgeschnittene Spitze (Grundfläche $9\,\mathrm{m}^2$, Höhe 6\,m); $V = 144 - 18 = 126\,\mathrm{m}^3$ \quad c) Erster Summand: Pyramide mit der Spitze auf der $z$-Achse bei 9 über dem rechtwinkligen Dreieck mit den Katheten 9 auf der $x$- und $y$-Achse; zweiter: die ähnliche Pyramide $OKLN$, deren Katheten und Höhe $k$ sind; $K$, $L$, $N$ liegen bei 4 auf den Achsen: $k = 4$ \quad d) $V_Q = G \cdot h$, $V_P = \tfrac{1}{3} \cdot G \cdot h$, $V_Q : V_P = 3$; Zuwachs $= 200\,\%$ \quad e) Kegel mit $r = 6$, $h = 8$, Mantellinie $s = 10$; $O = \pi \cdot 36 + \pi \cdot 6 \cdot 10 = 96\pi \approx 301{,}59$ \quad f) $10\,\mathrm{l} = 10\,000\,\mathrm{cm}^3$, $h = 10\,000 : (\pi \cdot 10^2) \approx 31{,}83\,\mathrm{cm}$ – höher als 30\,cm, sie passt nicht}
\erg{3}{a) $V(k) = \tfrac{1}{3} \cdot \tfrac{1}{2} \cdot (7 - k)(1 + k) \cdot 3 = \tfrac{1}{2} \cdot (7 - k)(1 + k)$: nach unten geöffnete Parabel mit den Nullstellen $-1$ und 7, das Maximum liegt in der Mitte: $k = 3$ \quad b) $V(k) = \tfrac{1}{3} \cdot \tfrac{1}{2} \cdot (4 + k)(8 - k) \cdot 9 = 1{,}5 \cdot (4 + k)(8 - k)$: nach unten geöffnete Parabel mit den Nullstellen $-4$ und 8, das Maximum liegt in der Mitte: $k = 2$ \quad c) $V_1 = \tfrac{1}{3} \cdot 36 \cdot (4 - k)$, $V_2 = \tfrac{1}{3} \cdot 36 \cdot (8 - (4 + k))$, zusammen $24 \cdot (4 - k) = 72$, $k = 1$ \quad d) $V_1 = \tfrac{1}{3} \cdot 20 \cdot (5 - k)$, $V_2 = \tfrac{1}{3} \cdot 20 \cdot (10 - (5 + k))$, zusammen $\tfrac{40}{3} \cdot (5 - k) = 40$, $k = 2$ \quad e) Graph II – $V(t) = 72 - \tfrac{1}{3} \cdot 12 \cdot (6 - 2t) = 48 + 8t$ ist linear in $t$, weil nur die Höhe der herausgenommenen Pyramide linear von $t$ abhängt; der Graph ist eine Gerade. \quad f) Graph I – $V(t) = 108 - \tfrac{1}{3} \cdot 12 \cdot (9 - 3t) = 72 + 12t$ ist linear in $t$, weil nur die Höhe der herausgenommenen Pyramide linear von $t$ abhängt; der Graph ist eine Gerade.}
